\clearpage

\section{Effective-spin description and Hamiltonian derivation}\label{app-nbcp-model-derivation}

\subsection{Magnetic degrees of freedom and the two-dimensional
approximation}\label{magnetic-degrees-of-freedom-and-the-two-dimensional-approximation}

The Co ions form triangular magnetic layers. Crystal-field effects and
atomic spin-orbit coupling produce a low-energy effective spin-\(1/2\)
doublet; higher multiplets lie approximately 71 meV above it. A dominant
in-plane nearest-neighbor model is motivated by the exchange paths and
weaker interlayer coupling. These are the material assumptions used in
Gao et al.~(2022)~\cite{gao2022},
rather than a claim that every further-neighbor or interlayer
interaction vanishes exactly.

Atomic spin-orbit coupling, which helps create the effective doublet,
must be distinguished from the coefficients of the residual
bond-dependent exchanges. A small fitted \(J_{\rm PD}\) or \(J_\Gamma\)
does not imply that the atomic spin-orbit interaction is absent. In this
note, ``zero SOC'' is shorthand for the exchange-model limit
\(J_{\rm PD}=J_\Gamma=0\).

We use \(z\parallel c\) and \(x\) along one nearest-neighbor bond. The
external field is \(B\parallel z\), and \(h=g_z\mu_B B\) is its Zeeman
energy. Spin operators are dimensionless, with eigenvalues measured in
units of \(\hbar\); exchange constants, \(h\), and excitation gaps are
expressed in meV. Energies denoted by \(e\) are per physical spin.

\subsection{Symmetry-constrained nearest-neighbor
exchange}\label{symmetry-constrained-nearest-neighbor-exchange}

The nearest-neighbor vectors in Fig.~\ref{fig-nbcp-bond-phases} are
\begin{equation*}
\boldsymbol\delta_1=a(1,0),\qquad
\boldsymbol\delta_2=a(-1/2,\sqrt3/2),\qquad
\boldsymbol\delta_3=a(-1/2,-\sqrt3/2),
\end{equation*}
where \(a\) is the lattice spacing. Translations along
\(\boldsymbol\delta_1\) and \(\boldsymbol\delta_2\) generate the lattice;
\(\boldsymbol\delta_3=-\boldsymbol\delta_1-\boldsymbol\delta_2\).
We first constrain the zero-field exchange tensor and then add the Zeeman
term. The symmetry argument fixes its form, not its coefficients~\cite{park2026}.

\subsubsection*{Bond-centered inversion.}
For one unoriented bond counted once, write
\(\hat H_{ij}=\hat{\mathbf S}_i^{\mathsf T}\mathsf J_{ij}\hat{\mathbf S}_j\).
Site-centered inversion followed by translation through
\(\boldsymbol\delta_1\) is inversion about the midpoint of the reference
bond. It exchanges \(i\) and \(j\). Spin is an axial vector, so
\(\mathcal I\hat{\mathbf S}_i\mathcal I^{-1}=\hat{\mathbf S}_j\)
without a component sign change. Since operators on different sites commute,
\begin{equation*}
\mathcal I\hat H_{ij}\mathcal I^{-1}
=\hat{\mathbf S}_j^{\mathsf T}\mathsf J_{ij}\hat{\mathbf S}_i
=\hat{\mathbf S}_i^{\mathsf T}\mathsf J_{ij}^{\mathsf T}\hat{\mathbf S}_j.
\end{equation*}
Invariance therefore requires \(\mathsf J_{ij}=\mathsf J_{ij}^{\mathsf T}\),
excluding antisymmetric exchange on this bond.

\subsubsection*{Twofold rotation about the reference bond.}
The rotation \(C_{2x}\) acts on the spin components through
\(D_x=\operatorname{diag}(1,-1,-1)\). The constraint
\(D_x^{\mathsf T}\mathsf J_0D_x=\mathsf J_0\) removes the \(xy\) and
\(xz\) entries, leaving
\begin{equation*}
\mathsf J_0=\begin{pmatrix}
J_{xx}&0&0\\0&J_{yy}&J_{yz}\\0&J_{yz}&J_{zz}
\end{pmatrix}.
\end{equation*}
The coefficients are defined by
\begin{equation*}
J=\frac{J_{xx}+J_{yy}}2,\quad J_z=J_{zz},\quad
J_{\rm PD}=\frac{J_{xx}-J_{yy}}4,\quad J_\Gamma=J_{yz}.
\end{equation*}
\newpage
These four independent coefficients give

\begin{equation}\phantomsection\label{eq-nbcp-reference-bond}{
\mathsf J_0=
\begin{pmatrix}
J+2J_{\rm PD}&0&0\\
0&J-2J_{\rm PD}&J_\Gamma\\
0&J_\Gamma&J_z
\end{pmatrix}.
}\end{equation}

\subsubsection*{Threefold rotation and the other bond directions.}
Use the active Cartesian rotation
\begin{equation*}
R_z(\beta)=\begin{pmatrix}
\cos\beta&-\sin\beta&0\\
\sin\beta&\cos\beta&0\\0&0&1
\end{pmatrix}.
\end{equation*}
Rotating the reference bond and both spins gives the covariance relation
\begin{align*}
\mathsf J_\beta&=R_z(\beta)\mathsf J_0R_z(\beta)^{\mathsf T}\\
&=\begin{pmatrix}
J+2J_{\rm PD}\cos2\beta&2J_{\rm PD}\sin2\beta&-J_\Gamma\sin\beta\\
2J_{\rm PD}\sin2\beta&J-2J_{\rm PD}\cos2\beta&J_\Gamma\cos\beta\\
-J_\Gamma\sin\beta&J_\Gamma\cos\beta&J_z
\end{pmatrix}.
\end{align*}
Only the discrete rotations \(\beta=0,\pm2\pi/3\) generate the three
bond representatives. Setting \(\beta=\varphi_d\), the identities
\(\cos2\varphi_d=\cos\varphi_d\) and
\(\sin2\varphi_d=-\sin\varphi_d\) give

\begin{equation}\phantomsection\label{eq-nbcp-bond-matrix}{
\mathsf J_d=
\begin{pmatrix}
J+2J_{\rm PD}\cos\varphi_d&-2J_{\rm PD}\sin\varphi_d&-J_\Gamma\sin\varphi_d\\
-2J_{\rm PD}\sin\varphi_d&J-2J_{\rm PD}\cos\varphi_d&J_\Gamma\cos\varphi_d\\
-J_\Gamma\sin\varphi_d&J_\Gamma\cos\varphi_d&J_z
\end{pmatrix}.
}\end{equation}

The microscopic Hamiltonian is consequently

\begin{equation}\phantomsection\label{eq-nbcp-model}{
\hat H=\sum_{\langle ij\rangle}\hat{\mathbf S}_i^{\mathsf T}
\mathsf J_{d(ij)}\hat{\mathbf S}_j-h\sum_i\hat S_i^z.
}\end{equation}

The bond phase \(\varphi_d\) labels a fixed exchange tensor. The
collective angle \(\phi\) introduced below rotates the spins while
keeping these tensors fixed. They must not be varied together when
evaluating the order-by-disorder energy.

With \(\hat S^\pm=\hat S^x\pm i\hat S^y\), the transverse XXZ
interaction is \(J(\hat S_i^+\hat S_j^-+\hat S_i^-\hat S_j^+)/2\). Thus
a convention with coefficient \(J_\pm\) multiplying that entire bracket
has \(J=2J_\pm\), not \(J=J_\pm\). The original notes contain both
identifications; the Cartesian matrix fixes the convention used here.

\newpage
\subsection{Parameter sets and the scope of the
calculation}\label{parameter-sets-and-the-scope-of-the-calculation}

The original notes contain several parameter sets with different
origins. They should remain separate when comparing a field quoted in
tesla with an exchange or field quoted in meV.

\begin{longtable}[]{@{}
  >{\raggedright\arraybackslash}p{(\linewidth - 6\tabcolsep) * \real{0.2500}}
  >{\raggedleft\arraybackslash}p{(\linewidth - 6\tabcolsep) * \real{0.2500}}
  >{\raggedleft\arraybackslash}p{(\linewidth - 6\tabcolsep) * \real{0.2500}}
  >{\raggedleft\arraybackslash}p{(\linewidth - 6\tabcolsep) * \real{0.2500}}@{}}
\toprule\noalign{}
\begin{minipage}[b]{\linewidth}\raggedright
Parameter set
\end{minipage} & \begin{minipage}[b]{\linewidth}\raggedleft
\(J\)
\end{minipage} & \begin{minipage}[b]{\linewidth}\raggedleft
\(J_z\)
\end{minipage} & \begin{minipage}[b]{\linewidth}\raggedleft
\(g_z\)
\end{minipage} \\
\midrule\noalign{}
\endfirsthead
\toprule\noalign{}
\begin{minipage}[b]{\linewidth}\raggedright
Parameter set
\end{minipage} & \begin{minipage}[b]{\linewidth}\raggedleft
\(J\)
\end{minipage} & \begin{minipage}[b]{\linewidth}\raggedleft
\(J_z\)
\end{minipage} & \begin{minipage}[b]{\linewidth}\raggedleft
\(g_z\)
\end{minipage} \\
\midrule\noalign{}
\endhead
\bottomrule\noalign{}
\caption{Exchange and Zeeman parameter sets and their distinct source conventions.}\label{tab-nbcp-02-model-hamiltonian-1}\\
\endlastfoot
Thermodynamic fit, Gao et al. & \(0.88\,k_B\mathrm{K}\) &
\(1.48\,k_B\mathrm{K}\) & 4.89 \\
Neutron fit, Woodland et al. & 0.0779 meV & 0.1225 meV & 4.716 \\
Older Overleaf phase calculations & 0.076 meV & 0.125 meV & 4.645 \\
Current Y/V gap comparison & 0.075 meV & 0.125 meV & 4.645 \\
\end{longtable}

The thermodynamic values are reported in
Gao et al.~\cite{gao2022}. The neutron values are Table~1 of
Woodland et al.~\cite{woodland2025}, whose \(J_{xy}\) multiplies
\(\hat S^x\hat S^x+\hat S^y\hat S^y\) and so equals \(J\) here; that fit
also gives \(g_{ab}=4.200\). It is the literature set of the toolkit
model definition. The last row follows the Figure~4 parameters of
Park et al.~\cite{park2026}. The middle row records the historical local
convention, not a new fit. The present calculation uses \(S=1/2\),
nearest-neighbor bonds and a single layer. The earlier
next-nearest-neighbor draft is not included: its axis and angular
conventions require further review.
